\documentclass[12pt, a4paper]{book}

% ─── Core packages ────────────────────────────────────────────────────────────
\usepackage[T1]{fontenc}
\usepackage[utf8]{inputenc}
\usepackage[charter, cal=cmcal]{mathdesign}
\usepackage[scaled=0.92]{helvet}
\usepackage{geometry}
\geometry{
  a4paper,
  top=30mm, bottom=30mm,
  inner=32mm, outer=25mm,
  marginparwidth=18mm
}

% ─── Mathematics ──────────────────────────────────────────────────────────────
\usepackage{amsmath, amsthm}

% ─── Chemistry ────────────────────────────────────────────────────────────────
\usepackage[version=4]{mhchem}

% ─── Graphics & colour ────────────────────────────────────────────────────────
\usepackage{graphicx}
\usepackage[dvipsnames,svgnames,table]{xcolor}
\usepackage{tikz}
\usetikzlibrary{arrows.meta, shapes.geometric, positioning,
                calc, backgrounds, decorations.pathmorphing,
                decorations.markings, patterns}
\usepackage{pgfplots}
\pgfplotsset{compat=1.18}
\usepackage{tcolorbox}
\tcbuselibrary{skins, breakable, theorems}

% ─── Layout ───────────────────────────────────────────────────────────────────
\usepackage{titlesec}
\usepackage{fancyhdr}
\usepackage{enumitem}
\usepackage{booktabs}
\usepackage{array}
\usepackage{multicol}
\usepackage{marginnote}
\usepackage{caption}
\usepackage{parskip}
\usepackage{siunitx}
\usepackage{pifont}
\usepackage[colorlinks=true, linkcolor=BrickRed,
            urlcolor=NavyBlue, citecolor=ForestGreen]{hyperref}

% ══════════════════════════════════════════════════════════════════════════════
% COLOUR PALETTE  (Journey 2 — deep purple)
% ══════════════════════════════════════════════════════════════════════════════
\definecolor{JourneyColor}{HTML}{4A235A}
\definecolor{ModuleColor}{HTML}{6C3483}
\definecolor{PuzzleColor}{HTML}{1A5276}
\definecolor{DiscoveryColor}{HTML}{B7950B}
\definecolor{WarningColor}{HTML}{C0392B}
\definecolor{ProblemColor}{HTML}{117A65}
\definecolor{ReflectColor}{HTML}{1D6A39}
\definecolor{HistoryColor}{HTML}{6E2F1A}
\definecolor{LightBg}{HTML}{F8F9FA}
\definecolor{ElectronBlue}{HTML}{1565C0}
\definecolor{NucleusRed}{HTML}{B71C1C}
\definecolor{AlphaGold}{HTML}{F57F17}

% ══════════════════════════════════════════════════════════════════════════════
% TCOLORBOX ENVIRONMENTS
% ══════════════════════════════════════════════════════════════════════════════
\newtcolorbox{puzzlebox}[1][]{
  enhanced, breakable,
  colback=PuzzleColor!6, colframe=PuzzleColor,
  fonttitle=\bfseries\sffamily, title={The Puzzle},
  attach boxed title to top left={yshift=-2mm, xshift=6mm},
  boxed title style={colback=PuzzleColor, colframe=PuzzleColor},
  borderline west={3pt}{0pt}{PuzzleColor}, #1}

\newtcolorbox{discoverybox}[1][]{
  enhanced, breakable,
  colback=DiscoveryColor!8, colframe=DiscoveryColor,
  fonttitle=\bfseries\sffamily, title={Discovery},
  attach boxed title to top left={yshift=-2mm, xshift=6mm},
  boxed title style={colback=DiscoveryColor, colframe=DiscoveryColor},
  borderline west={3pt}{0pt}{DiscoveryColor}, #1}

\newtcolorbox{problembox}[1][]{
  enhanced, breakable,
  colback=ProblemColor!5, colframe=ProblemColor,
  fonttitle=\bfseries\sffamily, title={Problems},
  attach boxed title to top left={yshift=-2mm, xshift=6mm},
  boxed title style={colback=ProblemColor, colframe=ProblemColor},
  borderline west={3pt}{0pt}{ProblemColor}, #1}

\newtcolorbox{deepbox}[1][]{
  enhanced, breakable,
  colback=JourneyColor!5, colframe=JourneyColor,
  fonttitle=\bfseries\sffamily, title={Deep Problems},
  attach boxed title to top left={yshift=-2mm, xshift=6mm},
  boxed title style={colback=JourneyColor, colframe=JourneyColor},
  borderline west={4pt}{0pt}{JourneyColor}, #1}

\newtcolorbox{historybox}[1][]{
  enhanced, breakable,
  colback=HistoryColor!7, colframe=HistoryColor!60,
  fonttitle=\bfseries\sffamily\small, title={Historical Note},
  attach boxed title to top left={yshift=-2mm, xshift=6mm},
  boxed title style={colback=HistoryColor!80, colframe=HistoryColor}, #1}

\newtcolorbox{reflectionbox}[1][]{
  enhanced, breakable,
  colback=ReflectColor!6, colframe=ReflectColor,
  fonttitle=\bfseries\sffamily, title={Reflection},
  attach boxed title to top left={yshift=-2mm, xshift=6mm},
  boxed title style={colback=ReflectColor, colframe=ReflectColor},
  borderline west={3pt}{0pt}{ReflectColor}, #1}

\newtcolorbox{modelbreaks}[1][]{
  enhanced, breakable,
  colback=WarningColor!6, colframe=WarningColor,
  fonttitle=\bfseries\sffamily, title={Where This Model Breaks},
  attach boxed title to top left={yshift=-2mm, xshift=6mm},
  boxed title style={colback=WarningColor, colframe=WarningColor},
  borderline west={4pt}{0pt}{WarningColor}, #1}

\newtcolorbox{namedconcept}[2][]{
  enhanced,
  colback=ModuleColor!6, colframe=ModuleColor,
  fonttitle=\bfseries\sffamily\large, title={#2},
  attach boxed title to top left={yshift=-2mm, xshift=6mm},
  boxed title style={colback=ModuleColor, colframe=ModuleColor}, #1}

% ══════════════════════════════════════════════════════════════════════════════
% THEOREM ENVIRONMENTS
% ══════════════════════════════════════════════════════════════════════════════
\theoremstyle{definition}
\newtheorem{observation}{Observation}[chapter]
\newtheorem{experiment}{Experiment}[chapter]

% ══════════════════════════════════════════════════════════════════════════════
% SECTION FORMATTING
% ══════════════════════════════════════════════════════════════════════════════
\titleformat{\chapter}[display]
  {\normalfont\sffamily\huge\bfseries\color{JourneyColor}}
  {\textcolor{JourneyColor!60}{\chaptertitlename\ \thechapter}}
  {10pt}{\Huge}[\vspace{2pt}\textcolor{JourneyColor}{\titlerule[1.5pt]}]

\titleformat{\section}
  {\normalfont\sffamily\Large\bfseries\color{ModuleColor}}
  {\thesection}{1em}{}[\textcolor{ModuleColor!40}{\titlerule[0.6pt]}]

\titleformat{\subsection}
  {\normalfont\sffamily\large\bfseries\color{JourneyColor!80}}
  {\thesubsection}{1em}{}

\titleformat{\subsubsection}
  {\normalfont\sffamily\normalsize\bfseries}
  {\thesubsubsection}{1em}{}

% ══════════════════════════════════════════════════════════════════════════════
% HEADER / FOOTER
% ══════════════════════════════════════════════════════════════════════════════
\pagestyle{fancy}
\fancyhf{}
\fancyhead[LE]{\small\sffamily\textcolor{JourneyColor}{\leftmark}}
\fancyhead[RO]{\small\sffamily\textcolor{ModuleColor}{\rightmark}}
\fancyfoot[C]{\small\thepage}
\renewcommand{\headrulewidth}{0.4pt}

% ══════════════════════════════════════════════════════════════════════════════
% CUSTOM COMMANDS
% ══════════════════════════════════════════════════════════════════════════════
\newcommand{\chem}[1]{\ce{#1}}
\newcommand{\important}[1]{\textcolor{ModuleColor}{\textbf{#1}}}
\newcommand{\newterm}[1]{\textit{\textcolor{DiscoveryColor}{#1}}}

% ══════════════════════════════════════════════════════════════════════════════
\begin{document}
\setcounter{chapter}{1}
\chapter[Journey 2 --- Architecture of the Atom]
        {Journey Two\\[8pt]
         \textcolor{ModuleColor}{\large Architecture of the Atom}}
\setcounter{section}{1}

% ══════════════════════════════════════════════════════════════════════════════
%  MODULE 2.2 — RUTHERFORD'S GOLD FOIL: THE NUCLEAR SURPRISE
% ══════════════════════════════════════════════════════════════════════════════
\section{Module 2.2 --- Rutherford's Gold Foil: The Nuclear Surprise}
\label{sec:mod2.2}

\subsection*{Motivation}

Thomson's plum pudding model made a clear, testable prediction:
fire small, fast, massive particles through an atom,
and they should pass straight through with only slight deflection.
The positive charge is so diffuse in the pudding model
that no single region of the atom can exert a large force.

Ernest Rutherford, who had studied under Thomson,
decided to test this prediction directly.
His tool: alpha particles emitted by radioactive materials.
His target: a thin sheet of gold foil.
His detector: a screen coated with zinc sulfide that flashes
when struck by an alpha particle.

What followed was one of the most dramatic experiments
in the history of science.

% ─── Section: Alpha Particles ─────────────────────────────────────────────────
\subsection{The Probe: Alpha Particles}

An \newterm{alpha particle} is the nucleus of a helium atom ---
two protons and two neutrons bound together.
Written as \ce{^4_2He^{2+}} or simply $\alpha$.

At the time of Rutherford's experiment (1909),
the nature of alpha particles was only partially understood.
What was known:
\begin{itemize}[leftmargin=2em, nosep]
  \item They are emitted by heavy radioactive atoms such as radium.
  \item They carry a positive charge (deflected by magnetic fields).
  \item They are relatively massive compared to electrons.
  \item They travel at roughly 5--7\% of the speed of light.
  \item They are stopped by a few centimetres of air or a sheet
        of paper.
\end{itemize}

Alpha particles were ideal probes precisely because
they are \emph{heavy} (about 8000 times the electron mass)
and \emph{fast}.
A massive, fast particle fired into a diffuse positive cloud
(the plum pudding) should barely notice the cloud is there.

\begin{center}
\begin{tikzpicture}[
  every node/.style={font=\sffamily\small, align=center},
  scale=0.95
]
  % Alpha particle properties
  \node[draw=AlphaGold!80, fill=AlphaGold!10, rounded corners=5pt,
        minimum width=12cm, minimum height=2.2cm]
        at (0,0) {};
  \node[font=\sffamily\small\bfseries, color=AlphaGold!80]
        at (0, 0.75) {Alpha Particle ($\alpha$)};
  \node[font=\sffamily\footnotesize, text=HistoryColor]
        at (-3.5, 0.05) {Charge: $+2e$\\$= +3.2\times10^{-19}$ C};
  \node[font=\sffamily\footnotesize, text=HistoryColor]
        at (0, 0.05) {Mass: $6.64\times10^{-27}$ kg\\$\approx 4$ u};
  \node[font=\sffamily\footnotesize, text=HistoryColor]
        at (3.5, 0.05) {Speed: $\approx 0.05c$\\$\approx 1.5\times10^7$ m/s};
  \node[font=\sffamily\footnotesize\itshape, text=JourneyColor!60]
        at (0,-0.65) {= helium nucleus; carries no electrons};
\end{tikzpicture}
\end{center}

% ─── Section: The Experimental Setup ──────────────────────────────────────────
\subsection{The Experimental Setup}

Rutherford and his students Hans Geiger and Ernest Marsden
set up the following apparatus:

\begin{center}
\begin{tikzpicture}[
  every node/.style={font=\sffamily\small},
  scale=1.0
]
  % Lead container with radioactive source
  \fill[JourneyColor!30] (-6.2,-0.9) rectangle (-5.2, 0.9);
  \fill[JourneyColor!60] (-6.2,-0.9) rectangle (-5.6, 0.9);
  \fill[AlphaGold!50] (-6.0,-0.2) circle (0.18cm);
  \node[font=\sffamily\footnotesize, text=HistoryColor]
    at (-5.7,-1.3) {Radium source in lead block};

  % Collimating slit
  \fill[JourneyColor!50] (-5.1,-0.8) rectangle (-4.8,-0.12);
  \fill[JourneyColor!50] (-5.1, 0.8) rectangle (-4.8, 0.12);
  \node[font=\sffamily\footnotesize, text=JourneyColor]
    at (-4.95,-1.3) {collimator};

  % Alpha beam
  \draw[-{Stealth[length=5pt]}, ultra thick, AlphaGold, line width=1.5pt]
    (-4.7, 0) -- (-1.5, 0);
  \node[font=\sffamily\footnotesize, text=AlphaGold]
    at (-3.2, 0.35) {$\alpha$ beam};

  % Gold foil
  \fill[DiscoveryColor!60] (-1.4,-1.5) rectangle (-1.2, 1.5);
  \node[font=\sffamily\footnotesize, text=DiscoveryColor]
    at (-1.3,-1.85) {gold foil (400 nm thick)};

  % Most go straight through
  \draw[-{Stealth[length=4pt]}, thick, AlphaGold!70]
    (-1.1, 0) -- (2.8, 0)
    node[above, font=\sffamily\footnotesize, text=AlphaGold!70]
      {most pass straight};

  % Small angle deflections
  \draw[-{Stealth[length=4pt]}, thick, AlphaGold!60]
    (-1.1, 0.05) -- (2.6, 0.6);
  \draw[-{Stealth[length=4pt]}, thick, AlphaGold!60]
    (-1.1,-0.05) -- (2.6,-0.6);

  % Large angle deflection
  \draw[-{Stealth[length=4pt]}, ultra thick, WarningColor]
    (-1.1, 0.05) -- (-2.5, 1.2)
    node[above left, font=\sffamily\footnotesize, text=WarningColor]
      {large angle back-scatter!};

  % ZnS detector screen (curved around)
  \draw[thick, ModuleColor!50, dashed]
    (3.5,-1.8) arc[start angle=-30, end angle=210, radius=3.5cm];
  \node[font=\sffamily\footnotesize, text=ModuleColor]
    at (2.5,-2.2) {ZnS detector screen};

  % Flashes
  \fill[DiscoveryColor!80] (3.0, 0) circle (0.08cm);
  \fill[DiscoveryColor!80] (2.8, 0.6) circle (0.06cm);
  \fill[DiscoveryColor!80] (2.8,-0.6) circle (0.06cm);
  \fill[WarningColor!80] (-3.1, 1.55) circle (0.08cm);
\end{tikzpicture}
\end{center}

\vspace{0.5em}

The zinc sulfide screen could be positioned anywhere around
the gold foil.
Geiger and Marsden patiently counted flashes,
sitting in the dark for hours as their eyes adjusted,
moving the screen to different angles.

% ─── Section: What They Expected ──────────────────────────────────────────────
\subsection{The Prediction vs the Reality}

\subsubsection{What the Plum Pudding Model Predicted}

If positive charge is uniformly distributed through the atom,
the maximum force on an incoming alpha particle at any point
is weak and points in no particular direction.
Simple calculations show:

\begin{center}
\begin{tcolorbox}[
  enhanced, colback=PuzzleColor!6, colframe=PuzzleColor!60,
  width=0.85\textwidth, arc=3pt
]
\centering\sffamily\small
\textbf{Plum pudding prediction:}\\[4pt]
Almost all alpha particles pass through with deflection less than 1\textdegree.\\
The probability of deflection greater than 90\textdegree\ is essentially zero ---\\
less than 1 in $10^{3500}$.
\end{tcolorbox}
\end{center}

\subsubsection{What They Actually Observed}

\begin{center}
\begin{tcolorbox}[
  enhanced, colback=WarningColor!8, colframe=WarningColor,
  width=0.85\textwidth, arc=3pt
]
\centering\sffamily\small
\textbf{Experimental result:}\\[4pt]
About 1 in 8000 alpha particles was deflected by \textbf{more than 90\textdegree}.\\
About 1 in 20000 bounced almost straight back (deflection $> 150$\textdegree).
\end{tcolorbox}
\end{center}

\vspace{0.5em}

Rutherford described his reaction with characteristic directness:

\begin{center}
\begin{tcolorbox}[
  enhanced, colback=HistoryColor!8, colframe=HistoryColor!50,
  width=0.85\textwidth, arc=3pt
]
\centering\sffamily\normalsize\itshape
``It was as if you fired a 15-inch artillery shell at a piece of tissue paper\\
and it came back and hit you.''\\[4pt]
\normalfont\small\sffamily --- Ernest Rutherford, 1913
\end{tcolorbox}
\end{center}

\vspace{0.5em}

For an alpha particle to bounce backward, it must encounter
a repulsive force far stronger than anything a diffuse
positive cloud could provide.
That force must come from something very small,
very dense, and very positively charged --- concentrated
in a tiny region, not spread through the whole atom.

\begin{puzzlebox}
Think through the physics carefully.

\begin{enumerate}[label=(\alph*), leftmargin=*]
  \item An alpha particle has charge $+2e$ and kinetic energy
        approximately $K = \SI{7.7}{\mega\electronvolt}$
        (about $1.23\times10^{-12}\ \si{\joule}$).
        For the alpha particle to be deflected backward,
        all its kinetic energy must be converted to
        electrostatic potential energy at closest approach.

        The electrostatic potential energy when two charges
        $q_1$ and $q_2$ are separated by distance $r$ is:
        \[
          U = \frac{q_1 q_2}{4\pi\varepsilon_0 r}
        \]
        where $\frac{1}{4\pi\varepsilon_0} = 8.99\times10^9\ \si{\newton\metre\squared\per\coulomb\squared}$.

        For gold ($Z = 79$), so $q_{\text{gold}} = 79e$.
        Set $K = U$ at closest approach and solve for $r$.

  \item Compare your answer from (a) to:
        \begin{itemize}[nosep]
          \item The radius of a gold atom: $\approx 1.44\times10^{-10}$ m
          \item The radius of the nucleus: $\approx 7.3\times10^{-15}$ m
        \end{itemize}
        What does this tell you about the size of the object
        that deflected the alpha particle?

  \item If the positive charge of gold were spread uniformly
        over a sphere of radius $1.44\times10^{-10}$ m (the atom),
        estimate the maximum force on an alpha particle
        passing through the surface.
        Compare this to the force at the nuclear radius.
        How many times stronger is the concentrated charge?

  \item Rutherford observed 1 in 20000 alphas bouncing back.
        This means 19999 in 20000 passed through.
        What does this ratio tell us about the
        \emph{fraction of the atom's volume} occupied
        by the nucleus?
\end{enumerate}
\end{puzzlebox}

% ─── Section: The Nuclear Model ───────────────────────────────────────────────
\subsection{Rutherford's Nuclear Model of the Atom}

From the scattering data, Rutherford derived a new model
of the atom (1911).
It was as radical as it was simple.

\begin{namedconcept}{The Nuclear Model of the Atom (Rutherford, 1911)}
\begin{itemize}[nosep, leftmargin=2em]
  \item The atom consists of a tiny, dense, positively charged
        \newterm{nucleus} at the centre.
  \item The nucleus contains almost all of the atom's mass.
  \item The nucleus is surrounded by electrons at relatively
        large distances.
  \item Most of the atom is \emph{empty space}.
\end{itemize}

\begin{center}
\begin{tabular}{lll}
\toprule
Property & Nucleus & Atom \\
\midrule
Radius & $\sim 10^{-15}\ \si{\metre}$ (1 fm) &
         $\sim 10^{-10}\ \si{\metre}$ (1 \AA) \\
Size ratio & \multicolumn{2}{l}{nucleus/atom $\approx 10^{-5}$
             (like a marble in a football stadium)} \\
Mass fraction & $> 99.97\%$ & 100\% (by definition) \\
Charge & $+Ze$ & neutral overall \\
\bottomrule
\end{tabular}
\end{center}
\end{namedconcept}

\vspace{0.5em}

The ``marble in a football stadium'' analogy deserves to be
felt, not just read.

\begin{center}
\begin{tikzpicture}[
  every node/.style={font=\sffamily\small, align=center},
  scale=0.85
]
  % Atom outline - large circle
  \draw[thick, JourneyColor!30, dashed]
    (0,0) circle (4.5cm);
  \node[font=\sffamily\footnotesize, text=JourneyColor!50]
    at (3.5, 3.8) {electron cloud\\radius $\sim 10^{-10}$ m};

  % Electrons scattered around
  \foreach \angle/\r in {
    30/3.8, 80/4.1, 140/3.9, 200/4.2,
    250/3.7, 310/4.0, 355/3.6, 170/4.3
  }{
    \fill[ElectronBlue] ({\r*cos(\angle)}, {\r*sin(\angle)}) circle (0.08cm);
  }

  % Nucleus - tiny dot at centre
  \fill[NucleusRed] (0,0) circle (0.14cm);
  \draw[thick, NucleusRed] (0,0) circle (0.14cm);
  \node[font=\sffamily\footnotesize, text=NucleusRed]
    at (1.1,-0.5) {nucleus\\$\sim 10^{-15}$ m\\(not to scale!)};
  \draw[->, thin, NucleusRed!60] (0.55,-0.35) -- (0.1,-0.05);

  % Alpha trajectories
  % Straight through (far from nucleus)
  \draw[-{Stealth[length=4pt]}, thick, AlphaGold!60]
    (-6.5, 3.2) -- (6.5, 3.2)
    node[right, font=\sffamily\tiny, text=AlphaGold!70]{undeflected};

  % Slight deflection
  \draw[-{Stealth[length=4pt]}, thick, AlphaGold!70]
    (-6.5, 1.0) .. controls (-2,0.8) and (2,0.6) .. (6.5, 0.4)
    node[right, font=\sffamily\tiny, text=AlphaGold!70]{slight deflect};

  % Large angle deflection (close to nucleus)
  \draw[-{Stealth[length=4pt]}, ultra thick, WarningColor]
    (-6.5,-0.3) .. controls (-1.5,-0.3) and (-0.5,0.5) ..
    (-1.5, 2.8)
    node[above, font=\sffamily\tiny, text=WarningColor]{large deflect};

  % Scale bar annotation
  \draw[{|}-{|}, thin, JourneyColor!40]
    (-4.5,-5.2) -- (4.5,-5.2)
    node[midway, below, font=\sffamily\footnotesize, text=JourneyColor!60]
      {atom diameter $\approx 2\times10^{-10}$ m};
\end{tikzpicture}
\end{center}

\vspace{0.5em}

This picture is startling.
An atom is almost entirely nothing.
If you scaled a hydrogen nucleus to the size of a marble (1 cm),
the electron would be orbiting at a distance of about 50 metres ---
half a football pitch away, in a stadium of empty space.

% ─── Section: The Rutherford Scattering Formula ───────────────────────────────
\subsection{The Rutherford Scattering Formula}

Rutherford did not just offer a qualitative picture.
He derived a precise mathematical formula for the fraction
of alpha particles deflected through any given angle $\theta$,
assuming all the positive charge is concentrated at a point:

\begin{namedconcept}{Rutherford Scattering Formula}
The fraction of alpha particles scattered through angles
\emph{greater than} $\theta$ is:

\[
  f(\theta) = \pi n t
  \left(\frac{Z_1 Z_2 e^2}{4K}\right)^2
  \cot^2\!\!\left(\frac{\theta}{2}\right)
\]

where:
\begin{itemize}[nosep, leftmargin=2em]
  \item $n$ = number density of atoms in the foil (atoms/m$^3$)
  \item $t$ = thickness of the foil
  \item $Z_1 = 2$ (alpha), $Z_2 = 79$ (gold)
  \item $K$ = kinetic energy of alpha particle
  \item $\theta$ = scattering angle
\end{itemize}

The key feature: this formula was \emph{derived from pure
Coulomb electrostatics} assuming a point nucleus.
Its agreement with experiment confirmed the nuclear model.
\end{namedconcept}

\vspace{0.5em}

The agreement between Rutherford's formula and Geiger and
Marsden's careful measurements was striking.
Every aspect of the angular distribution fit perfectly.
This was not an approximate agreement ---
it was a quantitative triumph.

\begin{historybox}
\textbf{Geiger, Marsden, and the Tedium of Discovery}

The gold foil experiment was largely carried out by Hans Geiger
and Ernest Marsden (who was an undergraduate student at the time).
They spent months in a darkened room, counting tiny flashes
on a zinc sulfide screen through a microscope, rotating
the detector to different angles, working in sessions that
lasted many hours.

Geiger later developed the \emph{Geiger counter} to automate
this kind of particle detection.
Marsden returned to New Zealand and became a physicist there.

The experiment took two years of careful work.
The theoretical analysis (Rutherford's formula) took
Rutherford about three months to derive.

This ratio --- years of careful data collection, months of
theoretical insight --- is typical of landmark experiments
in physics.
\end{historybox}

% ─── Section: Protons and Neutrons ────────────────────────────────────────────
\subsection{What Is the Nucleus Made Of?}

Rutherford's nuclear model established that almost all
the atom's mass is concentrated in the nucleus.
But what is the nucleus itself made of?

\subsubsection{The Proton}

Rutherford noticed in 1917 that when alpha particles
hit nitrogen gas, a new type of particle was ejected
with properties matching the hydrogen nucleus ---
a single positively charged particle with mass $\approx 1$ u.

\begin{center}
\ce{^{14}_{7}N + ^{4}_{2}He -> ^{17}_{8}O + ^{1}_{1}H}
\end{center}

He called this particle the \newterm{proton}.
Every nucleus contains protons.
The number of protons --- the \newterm{atomic number} $Z$ ---
determines which element the atom is.

\begin{namedconcept}{The Proton}
\begin{itemize}[nosep, leftmargin=2em]
  \item Charge: $+e = +1.602\times10^{-19}\ \si{\coulomb}$
  \item Mass: $m_p = 1.673\times10^{-27}\ \si{\kilogram} \approx 1$ u
  \item The number of protons $= Z =$ atomic number
  \item Determines the element's chemical identity
\end{itemize}
\end{namedconcept}

\subsubsection{The Neutron: A Neutral Companion}

There was a problem.
For all elements heavier than hydrogen,
the mass of the nucleus exceeds what the protons alone
could account for.
Helium has $Z = 2$ (2 protons) but mass $\approx 4$ u.
Where does the extra mass come from?

Rutherford predicted in 1920 that there must be a neutral
particle in the nucleus with approximately the mass of a proton.
He called it the \emph{neutron} --- but it took twelve years
to find.

In 1932, James Chadwick bombarded beryllium with alpha particles
and detected a neutral radiation that could knock protons
out of paraffin wax.
Only a particle with mass comparable to a proton could
do this; photons (light) could not.

\begin{center}
\ce{^{9}_{4}Be + ^{4}_{2}He -> ^{12}_{6}C + ^{1}_{0}n}
\end{center}

\begin{namedconcept}{The Neutron (Chadwick, 1932)}
\begin{itemize}[nosep, leftmargin=2em]
  \item Charge: 0 (neutral)
  \item Mass: $m_n = 1.675\times10^{-27}\ \si{\kilogram} \approx 1$ u
        (slightly heavier than proton)
  \item Number of neutrons $= N = A - Z$
        where $A$ is mass number
  \item Provides nuclear ``glue'' --- the strong nuclear force
        acts between neutrons and protons
\end{itemize}
\end{namedconcept}

\subsubsection{Nuclear Notation}

An atom is completely specified by two numbers:
$Z$ (protons, determines element) and
$A$ (mass number $=$ protons $+$ neutrons).

\[
  \ce{^{A}_{Z}X}
\]

For example:
\begin{itemize}[nosep, leftmargin=2em]
  \item \ce{^{1}_{1}H}: hydrogen, 1 proton, 0 neutrons
  \item \ce{^{4}_{2}He}: helium, 2 protons, 2 neutrons
  \item \ce{^{12}_{6}C}: carbon, 6 protons, 6 neutrons
  \item \ce{^{238}_{92}U}: uranium, 92 protons, 146 neutrons
\end{itemize}

\begin{center}
\begin{tikzpicture}[
  every node/.style={font=\sffamily\small, align=center},
  scale=1.0
]
  % Carbon-12 nucleus diagram
  \node[font=\sffamily\small\bfseries, color=JourneyColor]
    at (0, 2.8) {Carbon-12 nucleus (\ce{^{12}_{6}C})};

  % Draw protons and neutrons packed in nucleus
  \foreach \x/\y/\col/\lbl in {
    -0.4/0.4/NucleusRed/{$p^+$},
     0.4/0.4/NucleusRed/{$p^+$},
    -0.8/-0.1/NucleusRed/{$p^+$},
     0.0/-0.1/NucleusRed/{$p^+$},
     0.8/-0.1/NucleusRed/{$p^+$},
    -0.4/-0.6/NucleusRed/{$p^+$}
  }{
    \shade[ball color=\col!40] (\x,\y) circle (0.28cm);
  }
  \foreach \x/\y in {
    0.4/-0.6,
   -0.0/ 0.9,
    0.8/ 0.4,
   -0.8/ 0.9,
   -0.8/-0.6,
    0.4/-0.1
  }{
    \shade[ball color=JourneyColor!30] (\x,\y) circle (0.28cm);
  }

  % Legend
  \shade[ball color=NucleusRed!40] (2.4, 0.8) circle (0.22cm);
  \node[right, font=\sffamily\footnotesize, text=NucleusRed]
    at (2.65, 0.8) {proton ($p^+$)};
  \shade[ball color=JourneyColor!30] (2.4, 0.1) circle (0.22cm);
  \node[right, font=\sffamily\footnotesize, text=JourneyColor]
    at (2.65, 0.1) {neutron ($n^0$)};

  % Annotation
  \node[font=\sffamily\footnotesize, text=JourneyColor!70]
    at (0,-1.5) {6 protons + 6 neutrons = mass number 12};

  % Surrounding electron cloud hint
  \draw[thick, dashed, ElectronBlue!30]
    (0,0.15) circle (1.8cm);
  \foreach \angle in {0,60,120,180,240,300}{
    \fill[ElectronBlue] ({1.75*cos(\angle)},{0.15+1.75*sin(\angle)})
      circle (0.07cm);
  }
  \node[font=\sffamily\footnotesize, text=ElectronBlue!60]
    at (-2.4, 1.8) {6 electrons\\(far away)};
\end{tikzpicture}
\end{center}

% ─── Section: Isotopes Revisited ──────────────────────────────────────────────
\subsection{Isotopes: Same Element, Different Neutrons}

With the nuclear model, the concept of isotopes
becomes completely clear.
Isotopes of an element have the same number of protons
(same $Z$, same element) but different numbers of neutrons
(different $A$, different mass).

\begin{center}
\begin{tabular}{llccc}
\toprule
Element & Symbol & Protons ($Z$) & Neutrons ($N$) & Mass ($A$) \\
\midrule
Hydrogen    & \ce{^1H}   & 1 & 0 & 1 \\
Deuterium   & \ce{^2H}   & 1 & 1 & 2 \\
Tritium     & \ce{^3H}   & 1 & 2 & 3 \\
\midrule
Carbon-12   & \ce{^{12}C} & 6 & 6 & 12 \\
Carbon-13   & \ce{^{13}C} & 6 & 7 & 13 \\
Carbon-14   & \ce{^{14}C} & 6 & 8 & 14 \\
\midrule
Chlorine-35 & \ce{^{35}Cl}& 17 & 18 & 35 \\
Chlorine-37 & \ce{^{37}Cl}& 17 & 20 & 37 \\
\bottomrule
\end{tabular}
\end{center}

\vspace{0.5em}

Isotopes of the same element have \emph{identical chemical
properties} (because chemistry is determined by electrons,
which depends on proton number, not neutron number) but
\emph{different nuclear properties} (stability, radioactivity).

Carbon-14 is radioactive (unstable) because the neutron-to-proton
ratio is too high.
Carbon-12 and Carbon-13 are stable.

\begin{problembox}
\begin{enumerate}[label=\textbf{\arabic*.}]
  \item Complete the following table:
        \begin{center}
        \begin{tabular}{lccccc}
        \toprule
        Symbol & $Z$ & $N$ & $A$ & Element & Electrons (neutral atom) \\
        \midrule
        \ce{^{16}_{8}O}  & 8  & ?  & ?  & Oxygen  & ? \\
        \ce{^{56}_{26}Fe}& ?  & ?  & 56 & Iron    & ? \\
        \ce{^{238}_{92}U}& ?  & ?  & ?  & Uranium & ? \\
        \ce{^{131}_{53}I}& ?  & ?  & ?  & Iodine  & ? \\
        \bottomrule
        \end{tabular}
        \end{center}

  \item Natural chlorine is 75.77\% \ce{^{35}Cl} and 24.23\%
        \ce{^{37}Cl}.
        \begin{enumerate}[label=(\alph*)]
          \item Calculate the average atomic mass of chlorine.
          \item Explain why both isotopes have identical chemical
                behaviour but slightly different physical properties
                (e.g., different boiling points of \ce{HCl} made
                from each isotope).
        \end{enumerate}

  \item A radioactive decay series converts \ce{^{238}_{92}U}
        to \ce{^{206}_{82}Pb} through a sequence of alpha
        ($\alpha$) and beta ($\beta^-$) decays.
        An alpha decay: loses \ce{^4_2He} (loses 2 protons,
        2 neutrons).
        A beta decay: a neutron becomes a proton (gains 1 proton,
        loses 1 neutron).
        \begin{enumerate}[label=(\alph*)]
          \item How many protons and neutrons are lost overall
                going from \ce{^{238}_{92}U} to \ce{^{206}_{82}Pb}?
          \item How many alpha decays occurred?
          \item How many beta decays occurred?
        \end{enumerate}
\end{enumerate}
\end{problembox}

% ─── Section: The Size of the Nucleus ────────────────────────────────────────
\subsection{Quantifying the Nuclear Scale}

The nuclear model makes quantitative claims about sizes.
These can be tested.

The radius of a nucleus scales approximately as:
\[
  r = r_0 A^{1/3}
\]
where $r_0 \approx 1.2\times10^{-15}\ \si{\metre}$ (1.2 fm)
and $A$ is the mass number.

This ``$A^{1/3}$'' scaling is deeply significant:
it means the \emph{volume} of the nucleus is proportional
to $A$ (since $V \propto r^3 \propto A$).
Each nucleon (proton or neutron) occupies approximately
the same volume inside the nucleus.
Nuclear matter has roughly constant density.

\begin{namedconcept}{Nuclear Density}
The density of nuclear matter is approximately:
\[
  \rho_{\text{nuclear}} \approx 2 \times 10^{17}\ \si{\kilogram\per\metre\cubed}
\]
This is about $2 \times 10^{14}$ times denser than water.
A teaspoon of nuclear matter would weigh about
$10^{11}$ kg --- about 100 million tonnes.
\end{namedconcept}

\begin{problembox}
\begin{enumerate}[label=\textbf{\arabic*.}]
  \item Calculate the radius of the following nuclei
        using $r = r_0 A^{1/3}$, $r_0 = 1.2\times10^{-15}\ \si{\metre}$:
        \begin{enumerate}[label=(\alph*)]
          \item \ce{^{1}H} (hydrogen)
          \item \ce{^{4}He} (helium-4)
          \item \ce{^{12}C} (carbon-12)
          \item \ce{^{56}Fe} (iron-56)
          \item \ce{^{238}U} (uranium-238)
        \end{enumerate}

  \item For carbon-12, calculate:
        \begin{enumerate}[label=(\alph*)]
          \item The volume of the nucleus.
          \item The mass of the nucleus
                (use $m \approx A \times 1.67\times10^{-27}\ \si{kg}$).
          \item The density of the nucleus in \si{kg/m^3}.
          \item How does this compare to the density of water
                (\SI{1000}{kg/m^3})?
        \end{enumerate}

  \item The radius of a carbon atom is about
        $7\times10^{-11}$ m.
        \begin{enumerate}[label=(\alph*)]
          \item What fraction of the atomic radius is the nuclear radius?
          \item What fraction of the atomic \emph{volume} is the nuclear volume?
          \item If the atom were scaled up so the nucleus
                was the size of a tennis ball (radius 3.3 cm),
                how far away would the electrons be?
                Express in metres.
        \end{enumerate}
\end{enumerate}
\end{problembox}

% ─── Section: Where the Nuclear Model Breaks Down ─────────────────────────────
\subsection*{Where the Nuclear Model Breaks Down}

\begin{modelbreaks}
Rutherford's nuclear model solved the scattering problem
beautifully.
But it immediately raised a catastrophic new question.

If electrons are orbiting the nucleus like planets around
the sun, then by classical electrodynamics, they are
\emph{accelerating} (changing direction in circular motion).
Any accelerating charge radiates electromagnetic energy.
An orbiting electron should continuously lose energy,
spiral inward, and crash into the nucleus
in about $10^{-11}$ seconds.

But atoms are stable.
Hydrogen atoms have existed for billions of years.

The classical picture of orbiting electrons is therefore
\emph{wrong} --- but we do not yet know what to replace it with.
This is the crisis of classical physics.
Module~2.3 confronts it directly.

A second problem: if electrons orbit at any radius
(as in a classical planetary orbit), they should emit
a \emph{continuous} spectrum of light as they spiral in.
But excited atoms emit light only at specific, sharp
wavelengths --- a \emph{line spectrum}.

Both problems have the same solution: quantum mechanics.
But that solution requires completely abandoning the
idea of a classical trajectory.
\end{modelbreaks}

% ─── Section: Deep Problems ───────────────────────────────────────────────────
\subsection*{Deep Problems}

\begin{deepbox}
\begin{enumerate}[label=\textbf{D\arabic*.}]

  \item \textbf{Deriving the Distance of Closest Approach.}

        An alpha particle (charge $q_1 = +2e$, mass
        $m = 6.64\times10^{-27}$ kg) is fired directly at
        a gold nucleus (charge $q_2 = +79e$) with kinetic
        energy $K = \SI{7.7}{\mega\electronvolt}$.

        \begin{enumerate}[label=(\alph*)]
          \item Convert \SI{7.7}{\mega\electronvolt} to joules.
                ($\SI{1}{\electronvolt} = 1.602\times10^{-19}\ \si{\joule}$.)
          \item At the point of closest approach (direct
                head-on collision), all kinetic energy converts
                to electrostatic potential energy.
                Calculate the distance of closest approach $d$.
          \item Compare $d$ to the atomic radius of gold
                ($1.44\times10^{-10}$ m) and to the nuclear
                radius ($7.3\times10^{-15}$ m).
                Is the alpha particle probing the nucleus
                or just its vicinity?
          \item If the alpha particle had twice the kinetic energy,
                how would $d$ change?
                What does this tell us about using higher-energy
                probes to study nuclear structure?
        \end{enumerate}

  \item \textbf{The Strong Nuclear Force.}

        Protons in a nucleus are all positively charged
        and are packed within a radius of order $10^{-15}$ m.

        \begin{enumerate}[label=(\alph*)]
          \item Calculate the electrostatic repulsive force
                between two protons separated by
                $r = 2\times10^{-15}$ m.
          \item Convert this force to newtons.
                Comment on its magnitude.
          \item Protons clearly \emph{do} stay in the nucleus.
                This means there must be an attractive force
                stronger than the electrostatic repulsion at
                this range.
                This is the \emph{strong nuclear force}.
                From the force calculation in (a), what can
                you infer about the minimum strength of the
                strong nuclear force at nuclear distances?
          \item The strong nuclear force has a very short range ---
                it drops to essentially zero beyond about
                $2\times10^{-15}$ m.
                Explain qualitatively why this means that very
                large nuclei ($A > 200$) become unstable.
                (Hint: what happens to the electrostatic repulsion
                as the nucleus grows, while the strong force
                stays short-range?)
        \end{enumerate}

  \item \textbf{Rutherford Scattering and the Limit of the Model.}

        Rutherford's formula assumes the nucleus is a point charge.
        At low energies, this is a good approximation.
        At very high energies, alpha particles can actually
        reach the nuclear surface and the formula breaks down.

        \begin{enumerate}[label=(\alph*)]
          \item The nuclear radius of gold is approximately
                $r_{\text{Au}} = 7.3\times10^{-15}$ m.
                What kinetic energy must an alpha particle have
                to reach this distance from the gold nucleus?
                Express your answer in MeV.
          \item Real high-energy scattering experiments show
                deviations from Rutherford's formula at energies
                above the value you calculated.
                What does this deviation tell us?
          \item Rutherford's formula fits perfectly at low energies,
                implying the nucleus appears as a point charge.
                At high energies it fails.
                This is an example of a model having a domain
                of validity.
                Give two other examples from this curriculum
                of models with limited domains of validity.
        \end{enumerate}

\end{enumerate}
\end{deepbox}

% ─── Reflection ───────────────────────────────────────────────────────────────
\subsection*{Reflection}

\begin{reflectionbox}
\begin{enumerate}[label=\textbf{R\arabic*.}]
  \item Rutherford's result --- 1 in 20000 alphas bouncing back ---
        was completely unexpected.
        \textbf{What does this episode tell us about the
        importance of experimental anomalies in science?}
        How should a scientist respond when data contradicts
        a well-established model?

  \item The nuclear model says atoms are mostly empty space.
        \textbf{Why, then, does a table feel solid?}
        If your hand and the table are both mostly empty space,
        what is the physical origin of the force that prevents
        them from passing through each other?

  \item Thomson's plum pudding model was wrong,
        but it was wrong \emph{in a productive way} ---
        it made a precise prediction that Rutherford could test.
        \textbf{A model that makes no predictions cannot be wrong.}
        Is a non-falsifiable model useful in science?
        Give an example from this curriculum of a non-falsifiable
        or vague claim, and contrast it with a precise,
        testable one.

  \item We now have three subatomic particles:
        electrons, protons, and neutrons.
        \textbf{Which one determines:
        (a) the element's chemical identity;
        (b) the atom's mass;
        (c) whether the atom is electrically neutral;
        (d) whether the atom is radioactive?}
        Justify each answer.
\end{enumerate}
\end{reflectionbox}

\vspace{1.5em}
\begin{center}
\begin{tcolorbox}[
  enhanced, colback=DiscoveryColor!8, colframe=DiscoveryColor,
  width=0.88\textwidth, arc=3pt
]
\centering\sffamily\small
\textbf{What we invented in this module.}\\[4pt]
We invented the nucleus --- not by assuming it, but by firing
projectiles at atoms and watching where they went.
The gold foil experiment is a masterclass in
\emph{indirect inference}: you cannot see a nucleus,
but you can deduce its existence, size, and charge
from the pattern of deflections.
This technique --- scattering experiments to probe
invisible structure --- remains the foundation of
particle physics today.
The LHC at CERN is, in essence, a vastly more powerful
version of Rutherford's gold foil apparatus.\\[4pt]
But we now face a crisis.
The nuclear model demands that electrons orbit the nucleus ---
and classical physics says orbiting electrons must radiate
energy and spiral inward.
The atom should not exist.
That is the subject of Module~2.3.
\end{tcolorbox}
\end{center}

\vspace{2em}
\noindent\textcolor{JourneyColor!40}{\rule{\textwidth}{0.6pt}}
\begin{center}
\textit{\sffamily\small
Module 2.2 complete. Continue to Module 2.3:\\
\textbf{The Disaster of Classical Physics --- Why Atoms Shouldn't Exist}.}
\end{center}

\end{document}